% SPDX-License-Identifier: CC-BY-4.0
\section{Related Work}
\label{sec:related}

\subsection{Retrieval-Augmented Generation}

The dominant pattern for grounding LLM responses in external
knowledge couples a retrieval stage (chunking, embedding, top-$k$
nearest-neighbour search) with a generation stage (LLM
conditioned on the retrieved passages) \citep{lewis2020rag,
karpukhin2020dpr}. Substantial subsequent work improves the
retrieval side: HyDE generates a hypothetical answer with the LLM
and embeds that, recovering performance when the raw query is
short \citep{gao2022hyde}; cross-encoder reranking re-scores
candidates with $(query, passage)$ pairs scored directly
\citep{nogueira2019passage, reimers2019sentencebert};
RAG-Fusion generates multiple query rewrites and fuses retrieval
results with reciprocal rank fusion \citep{rackauckas2024ragfusion};
GraphRAG constructs an offline graph of the corpus and traverses
it at retrieval time \citep{edge2024graphrag}. Self-RAG trains
the LLM to adaptively invoke retrieval \citep{asai2024selfrag}.

These techniques operate on the retrieval pipeline; the generation
stage remains an LLM. We compare \hyphae{} against the vanilla
pipeline, HyDE plus cross-encoder reranking (the canonical
``serious RAG'' stack), and oracle retrieval (the LLM receives the
ground-truth supporting passage directly). The latter isolates
composition quality from retrieval quality.

\subsection{Attributed and Quote-Grounded Generation}
\label{sec:related-attribution}

A parallel line keeps the LLM in the answer path but makes the
output \emph{attributable}. Citation-generation systems emit an
answer with citation markers and are evaluated by frameworks such as
ALCE \citep{gao2023alce}, whose citation-NLI metric scores whether
the cited documents support the claim. \citet{menick2022gophercite}
(GopherCite) attaches \emph{verbatim} supporting quotes to generated
answers, learning the answer-to-quote support via reinforcement
learning from human preferences; ``according to'' prompting
\citep{weller2023accordingto} steers a model toward quoting its
pre-training data and measures the tendency with an n-gram-overlap
(QUIP) statistic. These establish verbatim quotation as an
attribution device as known prior art: \citet{schuster2023semqa}
(SEMQA) make it explicit, defining a \emph{semi-extractive} task
whose answers interleave verbatim source spans with generated
connectors for ``by-design'' easy-to-verify attribution.

\hyphae{} differs from all of these on one structural axis:
\emph{there is no model output to attribute}. The answer is not
generated and then linked to evidence; the answer \emph{is} the
stored fragment, emitted byte-for-byte, with the binding made
cryptographically tamper-evident rather than learned (GopherCite),
prompted (according-to), or textual-only (SEMQA, ALCE). This
collapses the distinction \citet{wallat2024faithfulness} draw
between citation \emph{correctness} (the cited source supports the
claim) and citation \emph{faithfulness} (the model genuinely relied
on the source rather than post-rationalising from its parametric
prior, a failure they find afflicts a large fraction of attributed
answers): a system with no parametric prior, whose answer \emph{is}
its source, is faithful by construction --- the post-rationalisation
failure mode is structurally eliminated, not merely reduced. It does
not collapse \emph{correctness}, which remains a retrieval-relevance
question; \citet{worledge2024extractive} characterise the resulting
extractive--abstractive verifiability trade-off, with the
byte-identical pole maximally verifiable and least abstractive. This
is consistent with our echo control: correctness is
realizer-independent, faithfulness is what verbatim emission makes
trivially true. Our provenance benchmark
(Section~\ref{sec:results-provenance}) is correspondingly orthogonal
to attribution benchmarks like ALCE --- it scores cryptographic
tamper \emph{detection and localisation} over the journal, not
citation correctness over the answer.

\subsection{Structured Memory and Cognitive Architectures}

Long-context architectures and explicit memory systems offer
different approaches to persistent context. MemGPT manages
context windows via a tiered storage system with LLM-mediated
paging \citep{packer2024memgpt}; Atlas (the production system
behind our \texttt{router:celiums-conversation} comparator) routes
queries across a curated model ensemble. Earlier work on cognitive
architectures (Soar \citep{laird2012cognitive}, ACT-R
\citep{anderson2014actr}) modelled memory and inference as
discrete operations on structured representations rather than as
prediction over token distributions.

\hyphae{} shares the structured-memory orientation but differs in
its composition stage: the realizer is not an LLM, not a learned
model, and not a statistical predictor. It is a finite procedure
walking a curated lexicon and emitting verbatim quotations of
retrieved fragments. The architecture is finite and inspectable;
no training run can change its outputs.

\subsection{Verifiable AI and NLI-Based Grounding}

Recent work on factual grounding evaluates LLM outputs against
their source contexts using natural language inference (NLI)
models \citep{honovich2022true, gekhman2023trueteacher}. The
\unsupf{} metric we employ is in this family: each sentence in
a response is scored for entailment against the retrieved
context, and sentences classified \texttt{neutral} or
\texttt{contradiction} count as unsupported. Recent surveys
\citep{huang2023hallucinationsurvey,
ji2023hallucinationsurvey} catalogue the failure modes such
metrics catch and miss.

We employ the \texttt{roberta-large-mnli} cross-encoder
\citep{liu2019roberta} for entailment scoring. Section~\ref{sec:setup}
documents the metric's known limitations
(asymmetric treatment of hedging, sensitivity to scaffolding
prose) and how \hyphae{}'s compositional template interacts
with them.

\subsection{Sub-millisecond Production Inference}

The latency-quality trade-off in conversational AI is widely
acknowledged, and a large line of work attacks LLM inference cost
directly --- e.g. speculative and staged-speculative decoding
\citep{leviathan2023speculative, chen2023fastinference}. Most
such work targets sub-second response under high token
generation budgets. \hyphae{} sits in a different operating
point: sub-millisecond composition under the structural
constraint that the response is bounded by curated lexicon
phrases and verbatim quotations. This shifts the operating
point of the system at the cost of generative flexibility
that some applications require and others do not.

\subsection{Tamper-Evident Logs and Provenance}

The cryptographic mechanism underlying \hyphae{}'s journal --- a
hash chain linking each record to its predecessor by digest --- is
classical, and we are explicit about that. Linked timestamping by
hash-chaining dates to \citet{haber1991timestamp}; Merkle trees
\citep{merkle1988digital} give the same tamper-evidence with
logarithmic verification; Certificate Transparency
\citep{laurie2014ct} deploys append-only Merkle logs at internet
scale with external auditing; and content-addressed DAGs such as
Git \citep{git} make the construction everyday infrastructure. We
claim none of this machinery as novel. What \hyphae{}'s journal
adds over a bare hash chain is only what every audited store needs
(durable persistence, an ordered chain over fragment ingests); the
contribution of this paper is not the chain.

Our contribution is an \emph{observation} about where such a chain
can and cannot be attached in a generation system: the audit
relation requires that the emitted output be byte-bindable to a
journalled source, and \textbf{paraphrastic generation destroys
that bindability} while \textbf{extractive (verbatim) emission
preserves it}. An LLM-RAG pipeline can hash-chain its corpus all it
likes; because its output paraphrases the retrieved text, no output
span is byte-identical to any journalled fragment, and the chain
cannot bind the answer the user received to the source it came
from. This is the property that divides auditable from
non-auditable retrieval systems, it is orthogonal to retrieval
quality and to the choice of hash construction, and to our
knowledge it has not been stated as the operative criterion for
provenance in retrieval-augmented generation.

\subsection{Where This Paper Fits}

We do not claim that \hyphae{} subsumes or replaces LLM-based
generation, and we do not claim it is more correct: our echo
control (Section~\ref{sec:results-triviaqa}) shows that on
standard correctness and grounding metrics a verbatim
\texttt{print} of the retrieved sentence matches it, so those
metrics measure verbatim quotation rather than the architecture.
We claim instead that \emph{for the subset of tasks in which the
unit of grounding is a discrete retrieved fragment and the
deployment requires provable provenance}, \hyphae{} provides a
property no LLM-RAG pipeline and no echo baseline provides by
construction: every emitted span is byte-identical to a fragment
in a hash-chained journal, so the output is
tamper-evident under the threat model of
Section~\ref{sec:results-provenance} --- auditable back to named,
unaltered source fragments against any attacker who does not hold
the external head-anchor key, which we implement and demonstrate. The trade-off, made
explicit, is that the prose is template-rigid; users wanting
LLM-style paraphrastic composition need an LLM.

The contribution is therefore positioned as a verifiable
\emph{provenance layer} for grounded retrieval, not a
correctness-competitive replacement for LLM-RAG.
